\documentclass[12pt,a4paper]{exam}
\usepackage[T1]{fontenc}
\usepackage{longtable}
\usepackage{graphicx}

\usepackage{hyperref}
%\urlstyle{same}  % don't use monospace font for urls
\setlength{\parindent}{0pt}
\setlength{\parskip}{6pt plus 2pt minus 1pt}

\usepackage{geometry}
\usepackage{multicol}
\usepackage{relsize}
\usepackage{pgf}
\usepackage{minted}

\geometry{
  %includeheadfoot,
  margin=2.54cm
}
\title{Exercises on Design Patterns III}
\author{Some more patterns}
\date{2018}

\fvset{tabsize=2}
%\fvset{fontsize=\relsize{-1}}
%\fvset{frame=single}

% https://www.javacodegeeks.com/2015/09/decorator-design-pattern.html

\newcommand{\code}{/Users/keith/Documents/Courses/sdp/repo/worksheets/dp-three/kotlin/src}
%\newcommand{\codesol}{/Users/keith/Courses/sdp/SDP-2017/exercises/week10sol/src/main/scala}

\begin{document}
\maketitle

In these exercises we will be examining the following design patterns:
\begin{itemize}
\item Chain of Responsibility,
\item Command,
\item Interpreter,
\item Prototype,
\item Iterator, and
\item Mediator.
\end{itemize}

The source code snippets listed below are written using the Kotlin language; similar versions in Scala can be 
found on the source code repository under the \verb!worksheets/dp-three! folder. 

\begin{questions}

\question
This question concerns the \textsc{Chain of Responsibility} design pattern.

Your company has got a contract to provide an analytical application to a health company. The application would 
tell the user about the particular health problem, its history, its treatment, medicines, 
interview of the person suffering from it, etc.; basically everything that is required to know about it. 
For this, your company receives a huge amount of data. The data could be in any format, it could text files, doc 
files, excels, audio, images, videos, anything that you can think of would be there.

Now, your job is to save this data in the company's database. Users will provide the data in any format and you 
should provide them a single interface to upload the data into the database. The user is not interested, not even 
aware, to know that how you are saving the different unstructured data?

The problem here is that you need to develop different handlers to save the various formats of data. For example, 
a text file save handler does not know how to save an mp3 file.

To solve this problem you can use the \emph{Chain of Responsibility} design pattern. You can create different 
objects which process different formats of data and chain them together. 
When a request comes to a single object, it will check whether it can process and handle the specific file format. 
If it can, it will process it; otherwise, it will forward it to the next object chained to it. 
This design pattern also decouples the user from the object that is serving the request; the user is not aware 
which object is actually serving its request.

To implement the Chain of Responsibility pattern to solve the above problem, you will create an interface, 
\mintinline{scala}{Handler}.
\VerbatimInput{\code/chain/Handler.kt}
The above interface contains two main methods, the \mintinline{scala}{setHandler} and the \mintinline{scala}
{process} methods. The \mintinline{scala}{setHandler} is used to set the next handler in the 
chain, whereas; the \mintinline{scala}{process} method is used to process the request, only if the handler can 
able process the request. 
Optionally, we have the \mintinline{scala}{getHandlerName} method which is used to return the handler's name.
The handlers are designed to process files which contain data. The concrete handler checks if it's able to handle 
the file by checking the file type, otherwise forwards to the next handler in the chain.

The \mintinline{scala}{File} class looks like this:
\VerbatimInput{\code/chain/File.kt}
The \mintinline{scala}{File} class creates simple file objects which contain the file name, file type, and the 
file path. 
The file type would be used by the handler to check if the file can be handled by them or not. If a handler can, 
it will process and save it, or it will forward it to the next handler in the chain.

Your solution should satisfy the following tests:
\VerbatimInput{\code/chain/TestChainofResponsibility.kt}
The above program should have the following output:
\begin{Verbatim}
Text Handler forwards request to Doc Handler
Doc Handler forwards request to Excel Handler
Excel Handler forwards request to Audio Handler
Process and saving audio file... by Audio Handler

Text Handler forwards request to Doc Handler
Doc Handler forwards request to Excel Handler
Excel Handler forwards request to Audio Handler
Audio Handler forwards request to Video Handler
Process and saving video file... by Video Handler

Text Handler forwards request to Doc Handler
Process and saving doc file... by Doc Handler

Text Handler forwards request to Doc Handler
Doc Handler forwards request to Excel Handler
Excel Handler forwards request to Audio Handler
Audio Handler forwards request to Video Handler
Video Handler forwards request to Image Handler
File not supported
\end{Verbatim}
\begin{solution}
\ldots	
\end{solution}


\question
This question concerns the \textsc{Command} design pattern.

You are required to create an application to execute different types of jobs. A job can be anything in a system,
for example, sending emails, SMS, logging, or performing some IO functions.
The \textsc{Command} pattern will help to decouple the invoker from a receiver and helps to execute any type of 
job without knowing its implementation. 
We'll make this example more interesting by creating threads which will help to execute these jobs concurrently. 
As these jobs are independent of each other, so the sequence of execution of these 
jobs is not really important. We will create a thread pool to limit the number of threads to execute jobs. A 
command object will encapsulate jobs and will hand it over to a thread from the pool 
that will execute the job.

The command object will be referenced by a common interface and will contain a method that will be used to execute 
the requests. 
The concrete command classes will override that method and will provide their own specific implementation to 
execute the request.
\VerbatimInput{\code/command/Job.kt}
The \mintinline{scala}{Job} interface is the command interface, contains a single method \mintinline{scala}{run}, 
which is executed by a thread. 
Our command's \mintinline{scala}{execute} method is the \mintinline{scala}{run} method which will be used to 
execute by a thread to get the work done.
There would be a different type of jobs that can be executed. The following is one of the concrete classes whose 
instances will be executed by the different command objects:
\VerbatimInput{\code/command/Email.kt}

The implementation classes of the \mintinline{scala}{Job} interface hold a reference to their respective classes 
that will be used to get the job done. 
The classes override the \mintinline{scala}{run} method and do the work requested. For example, the 
\mintinline{scala}{SmsJob} class is used to send sms, its 
\mintinline{scala}{run} method calls the \mintinline{scala}{sendSms} method of the \mintinline{scala}{Sms} object 
to get the job done.

You can set different objects one by one to the same command object.

The below is the \mintinline{scala}{ThreadPool} class used to create pool of threads and allow a thread to fetch 
and execute the job from the job queue.
\VerbatimInput{\code/command/ThreadPool.kt}
The above class is used to create $n$ threads (worker threads). Each worker thread will wait for a job in a queue 
and then execute the job and will go back to waiting state. 
The class contains a job queue; when a new job will be added into the queue, a worker thread from the pool will 
execute the job.

Your solution should satisfy the following tests:
\VerbatimInput{\code/command/TestCommandPattern.kt}
The above code will result in the following output:
\begin{Verbatim}
Job ID: 9 executing email jobs.
Sending email.......
Job ID: 12 executing logging jobs.
Job ID: 17 executing email jobs.
Sending email.......
Job ID: 13 executing email jobs.
Sending email.......
Job ID: 10 executing sms jobs.
Sending SMS...
Job ID: 11 executing fileIO jobs.
Executing File IO operations...
Job ID: 18 executing sms jobs.
Sending SMS...
Logging...
Job ID: 16 executing logging jobs.
Logging...
Job ID: 15 executing fileIO jobs.
Executing File IO operations...
Job ID: 14 executing sms jobs.
Sending SMS...
Job ID: 12 executing fileIO jobs.
Executing File IO operations...
Job ID: 10 executing logging jobs.
Logging...
Job ID: 18 executing email jobs.
Sending email.......
Job ID: 16 executing sms jobs.
Sending SMS...
Job ID: 14 executing fileIO jobs.
Executing File IO operations...
Job ID: 9 executing logging jobs.
Logging...
Job ID: 17 executing email jobs.
Sending email.......
Job ID: 13 executing sms jobs.
Sending SMS...
Job ID: 15 executing fileIO jobs.
Executing File IO operations...
Job ID: 11 executing logging jobs.
Logging...
\end{Verbatim}
Due to the nature of concurrency the output may differ on subsequent executions.
\begin{solution}
	
\end{solution}

\question
This question concerns the \textsc{Interpreter} design pattern.

Given a language, define a representation for its grammar along with an interpreter that uses the 
representation to interpret sentences in the language.

In general, languages are made up of a set of grammar rules. Different sentences can be constructed by 
following these grammar rules. Sometimes an application may need to process repeated occurrences of 
similar requests that are a combination of a set of grammar rules. These requests are distinct but are 
similar in the sense that they are all composed using the same set of rules.

A simple example of this would be the set of different arithmetic expressions submitted to a calculator program 
(some familiar?). Though each such expression is different, they are all constructed using the basic rules 
that make up the grammar for the language of arithmetic expressions.

In such cases, instead of treating every distinct combination of rules as a separate case, it may be beneficial 
for the application to have the ability to interpret a generic combination of rules. The Interpreter pattern can 
be used to design this ability in an application so that other applications and users can specify operations using 
a simple language defined by a set of grammar rules.

A class hierarchy can be designed to represent the set of grammar rules with every class in the hierarchy 
representing a separate grammar rule. An Interpreter module can be designed to interpret the sentences constructed 
using the class hierarchy designed above and carries out the necessary operations.

Because a different class represents every grammar rule, the number of classes increases with the number of 
grammar rules. A language with extensive, complex grammar rules requires a large number of classes. The 
\textsc{Interpreter} pattern works best when the grammar is simple. Having a simple grammar avoids the need 
to have many classes corresponding to the complex set of rules involved, which are hard to manage and maintain.

Consider the following interface specification:
\VerbatimInput{\code/interpreter/Expression.kt}
The above interface is used by all different concrete expressions and overrides the interpret method to 
define their specific operation on the expression.
The following are the operation specific expression classes which you need to complete:
\VerbatimInput{\code/interpreter/Add.kt}
\VerbatimInput{\code/interpreter/Product.kt}
Your solution should satisfy the following tests:
\VerbatimInput{\code/interpreter/TestInterpreterPattern.kt}
The code above should provide the following output:
\begin{Verbatim}
( 7 3 -2 1 + * ):12
\end{Verbatim}
Please note that we have used postfix notation.

\subsubsection*{Postfix notation}

If you don't know about postfix, here is a brief introduction. There are three notations for 
a mathematical expression i.e., \emph{infix}, \emph{postfix}, and \emph{prefix}.
\begin{description}
\item[Infix notation] is the common arithmetic and logical formula notation, in which operators are written infix-
style between the operands they act on e.g. $3 + 4$.
\item[Postfix a.k.a. Reverse Polish notation (RPN)] is mathematical notation in which every operator follows all 
of its operands e.g. $3 4 +$.
\item[Prefix (Polish notation)] is a form of notation for logic, arithmetic, and algebra in which operators to the 
left of their operands e.g. $+ 3 4$.
\end{description}
Infix notation is a normally used in mathematical expressions. The other two notations are used as syntax for 
mathematical expressions by interpreters of programming languages.

In the above class, we declared a postfix of an expression in the \mintinline{scala}{tokenString} variable. 
Then we split the \mintinline{scala}{tokenString} and assigned it into an array, the \mintinline{scala}
{tokenArray}. While iterating tokens one by one, first we have checked whether the token 
is an \emph{operator} or an \emph{operand}. If the token is an operand we pushed it onto the stack, but if it is 
an operator we popped the first two operands from the stack. 
The \mintinline{scala}{getOperation} method returns the appropriate expression class according to the operator 
passed to it.
We next interpreted the result and pushed it back onto the stack. After iterating over the complete 
\mintinline{scala}{tokenList} we got the final result.

\begin{solution}
	
\end{solution}

\question
This question concerns the \textsc{Prototype} design pattern.

Let us consider a scenario where an application requires some access control. The features of the applications 
can be used by the users according to the access rights provided to them. For example, some users have access to 
the reports generated by the application, while some don't. Some of them even can modify the reports, while some 
can only read it. Some users also have administrative rights to add or even remove other users.

Every user object has an access control object, which is used to provide or restrict the controls of the 
application. This access control object is a bulky, heavy object and its creation is very costly since it requires 
data to be fetched from some external resources, like databases or some property files etc.

We also cannot share the same access control object with users of the same level, because the rights can be 
changed at runtime by the administrator and a different user with the same level could have a different access 
control. One user object should have one access control object.

We can use the Prototype Design Pattern to resolve this problem by creating the access control objects on all 
levels at once, and then provide a copy of the object to the user whenever required. In this case, data fetching 
from the external resources happens only once. Next time, the access control object is created by copying the 
existing one. The access control object is not created from scratch every time the request is sent; this approach 
will certainly reduce object creation time.

You will use the \mintinline{scala}!clone! method to solve the above problem.
\VerbatimInput{\code/prototype/Prototype.kt}
The above interface extends the \mintinline{scala}!Cloneable! interface and contains a method 
\mintinline{scala}!clone!. This interface is implemented by classes which want to create a prototype object.

The \mintinline{scala}!AccessControl! class implements the 
\mintinline{scala}!Prototype! interface and overrides the \mintinline{scala}!clone! method.
The class also contains two properties; the \mintinline{scala}!controlLevel! is used to 
specific the level of control this object contains. The level depends upon the type of user going to use it, 
for example, USER, ADMIN, MANAGER etc.

The other property is the \mintinline{scala}!access!; it contains the access right for the user. 
Please note that, for simplicity, we have used access as a \mintinline{scala}!String! type attribute. 
This could be of type \mintinline{scala}!Map! which can contain key value pairs of long access 
rights assigned to the user.

The \mintinline{scala}!User! class has a \mintinline{scala}!userName!, 
\mintinline{scala}!level! and a reference to the \mintinline{scala}!AccessControl! assigned to it.

You should use an \mintinline{scala}!AccessControlProvider! class that creates and 
stores the possible \mintinline{scala}!AccessControl! objects in advance.
When there is a request to an \mintinline{scala}!AccessControl! object, 
it returns a new object created by copying the stored prototypes.

The \mintinline{scala}!getAccessControlObject! method fetches a stored prototype 
object according to the \mintinline{scala}!controlLevel! passed to it, 
from the map and returns a newly created cloned object to the client code.

You should test your code using the following:
\VerbatimInput{\code/prototype/TestPrototypePattern.kt}
The code above should provide the following output:
\begin{Verbatim}
Fetching data from external resources and creating access control objects...
************************************
Name: User A, Level: USER Level, Access Control Level:USER, Access: DO_WORK
Changing access of: User B
Name: User B, Level: USER Level, Access Control Level:USER, Access: READ REPORTS 
************************************
Name: User C, Level: MANAGER Level, Access Control Level:MANAGER, Access: 
	GENERATE/READ REPORTS	
\end{Verbatim}

\begin{solution}
	
\end{solution}

\question
This question concerns the \textsc{Iterator} design pattern which you are going to explore 
using a \mintinline{scala}!Shape! class. We will store and iterate through the \mintinline{scala}!Shape! 
objects using an iterator.
\VerbatimInput{\code/iterator/Shape.kt}
The simple \mintinline{scala}!Shape! class has an \mintinline{scala}!id! and 
\mintinline{scala}!name! as its attributes.
\VerbatimInput{\code/iterator/ShapeStorage.kt}
The above class is used to store the \mintinline{scala}!Shape! objects. 
The class contains an array of \mintinline{scala}!Shape! type; for simplicity we have initialised 
that array up to five elements. The \mintinline{scala}!addShape! method is used to add a 
\mintinline{scala}!Shape! object to the array and increment the index by one. The \mintinline{scala}!getShapes!
method returns the array of \mintinline{scala}!Shape! type.

The \mintinline{scala}!ShapeIterator! class is an \mintinline{scala}!Iterator! to the 
\mintinline{scala}!Shape! class. The class implements the \mintinline{scala}!Iterator! 
interface and defines all the methods of the \mintinline{scala}!Iterator! interface.

The \mintinline{scala}!hasNext! method returns a \mintinline{scala}!Boolean! if there's an item left. 
The \mintinline{scala}!next! method returns the next item from the collection and the 
\mintinline{scala}!remove! method removes the current item from the collection.

You should test your code using the following:
\VerbatimInput{\code/iterator/TestIteratorPattern.kt}
The above code should result to the following output:
\begin{Verbatim}
ID: 0 Shape: Polygon
ID: 1 Shape: Hexagon
ID: 2 Shape: Circle
ID: 3 Shape: Rectangle
ID: 4 Shape: Square
Apply removing while iterating...
ID: 0 Shape: Polygon
ID: 2 Shape: Circle
ID: 4 Shape: Square	
\end{Verbatim}

\begin{solution}
	
\end{solution}

\question
This question concerns the \textsc{Mediator} design pattern.

A big electronic company has asked you to develop a piece of software to operate its new fully automatic 
washing machine. The company has provided you with the hardware specification and the working knowledge of 
the machine. In the specification, they have provided you the different washing programs the machine 
supports. They want to produce a fully automatic washing machine that will require almost zero human 
intervention. 
The user should only need to connect the machine with a tap to supply water, load the clothes to wash, 
set the type of clothes in the machine like cotton, silk, or denims etc, provide detergent and softener 
to their respective trays, and press the start button.

The machine should be smart enough to fill the water in the drum, as much as required. It should 
adjust the wash temperature by itself by turning the heater on, according to the type of clothes in it. 
It should start the motor and spin the drum as much required, rinse according to the clothes needs, use 
soil removal if required, and softener too.

As an Object Oriented developer, you started analysing and classifying objects, classes, and their relationships. 
Let's check some of the important classes and parts of the system. First of all, a 
\mintinline{scala}!Machine! class, which has a drum. So a 
\mintinline{scala}!Drum! class, but also a heater, a sensor to check the temperature, a motor. 
Additionally, the machine also has a valve to control the water supply, a soil removal, detergent, and a softener.

These classes have a very complex relationship with each other, and the relationship also varies. 
Please note that currently we are talking only about the high level abstraction of the machine. 
If we try to design it without keeping much of OOP principles and patterns in mind, then the initial 
design would be very tightly coupled and hard to maintain. 
This is because the above classes should contact with each other to get the job done. For example, 
the \mintinline{scala}!Machine! class should ask the \mintinline{scala}!Valve! class to open the valve, 
or the \mintinline{scala}!Motor! should spin the 
\mintinline{scala}!Drum! at certain rpm according to the wash program set (which is set by the type of clothes 
in the machine). Some type of clothes require softener or soil removal while others don't, or the temperature
should be set according to the type of clothes.

If we allow classes to contact each other directly, that is, by providing a reference, 
the design will become very tightly coupled and hard to maintain. 
It would become very difficult to change one class without affecting the other. 
Even worse, the relationship between the classes varies, according to the different wash programs, 
like different temperature for different type of clothes. So these classes won't able to be reused. 
To support all the wash programs we need to put control statements like \emph{if-else} 
in the code which would make the code even more complex and difficult to maintain.

To decouple these objects from each other we need a \textsc{Mediator}, which will contact the object on 
behalf of the other object, hence providing loose coupling between them. 
The object only needs to know about the mediator, and perform operations on it. 
The mediator will perform operations on the required underlying object in order to get the work done.

You will see how the \textsc{Mediator Pattern} will make the washing machine design better, 
reusable, maintainable and loosely coupled.
\VerbatimInput{\code/mediator/MachineMediator.kt}
The \mintinline{scala}!MachineMediator! is an interface which acts as a generic mediator. 
\VerbatimInput{\code/mediator/Colleague.kt}
The \mintinline{scala}!Colleague! interface has one method to set the mediator for the concrete 
colleague's class.
\VerbatimInput{\code/mediator/Button.kt}
The above \mintinline{scala}!Button! class is a colleague class which holds a reference to a mediator. 
The user presses the button which calls the \mintinline{scala}!press! method of this class which in turn, 
calls the \mintinline{scala}!start! method of the concrete mediator class. 
This \mintinline{scala}!start! method of the mediator calls the 
\mintinline{scala}!start! method of the machine class on behalf of the \mintinline{scala}!Button! class.

As stated by the company, the washing machine has a set of wash programs and the software should support 
all these programs. 
The following mediator is set as a washing program for cottons, so parameters such as temperature, 
drum spinning speed, level of soil removal etc., are set accordingly. 
We can have different mediators for different washing programs without changing the existing colleague 
classes and thus providing loose coupling and reusability. 
All these colleague classes can be reused with others washing programs of the machine.
\VerbatimInput{\code/mediator/CottonMediator.kt}

To see the advantages and power of the \textsc{Mediator Pattern}, let's take another mediator that is 
used as a washing program for denim jeans. Now we just need to create a new mediator and set the rules in 
it to wash denims: rules like temperature, and the speed at which drum will spin, whether softener is 
required or not, the level of the soil removal, etc. We don't need to change anything in the existing structure.
No conditional statements like \emph{if-else} are required, something that would otherwise 
increase the complexity.
\VerbatimInput{\code/mediator/DenimMediator.kt}
Now we test these mediators:
\VerbatimInput{\code/mediator/TestMediator.kt}
and should result in the following output:
\begin{Verbatim}
Setting up for COTTON program
Button pressed.
Valve is opened...
Filling water...
Valve is closed...
Heater is on...
Temperature reached 40 C
Temperature is set to 40
Heater is off...
Start motor...
Rotating drum at 700 rpm.
Adding detergent
Setting Soil Removal to low
Adding softener

Setting up for DENIM program
Button pressed.
Valve is opened...
Filling water...
Valve is closed...
Heater is on...
Temperature reached 30 C
Temperature is set to 30
Heater is off...
Start motor...
Rotating drum at 1400 rpm.
Adding detergent
Setting Soil Removal to medium
No softener is required	
\end{Verbatim}

\begin{solution}
	
\end{solution}
\end{questions}

\end{document}